\chapter{Mack's lognormal approximation}

The parameter conversion appears in Thomas Mack, ``Measuring the Variability
of Chain Ladder Reserve Estimates,'' Casualty Actuarial Society Forum, Spring
1994, Chapter 4, equation (10), pages 118--119. The 1993 ASTIN paper does not
contain this construction.

\begin{definition}[Moment-matched parameters]\label{def:mack-lognormal}\lean{VerifiedReserving.mackLognormalSigma2, VerifiedReserving.mackLognormalMu}\leanok
For an arithmetic mean $m$ and variance $v$, define
\[
  \sigma^2=\log(1+v/m^2),
  \qquad
  \mu=\log m-\sigma^2/2.
\]
These are the log-variance and log-mean parameters in Mack's equation (10).
\end{definition}

\begin{theorem}[Matched lognormal moments]\label{thm:mack-lognormal}\lean{VerifiedReserving.mackLognormal_mean_eq, VerifiedReserving.mackLognormal_variance_eq}\leanok
\uses{def:mack-lognormal}
If $m>0$ and $v\ge0$, then
\[
  \exp(\mu+\sigma^2/2)=m
\]
and
\[
  \exp(2\mu+\sigma^2)\bigl(\exp(\sigma^2)-1\bigr)=v.
\]
Thus the standard lognormal moment expressions reproduce the supplied mean
and variance. This identity does not state that the underlying reserve follows
a lognormal law.
\end{theorem}
\begin{proof}\leanok Substitute the parameter definitions, use
$\exp(\log x)=x$ for positive $x$, and simplify the resulting rational
identity.
\end{proof}
